\section{Implementation}
\label{sec:implementation}

The prototype uses three Kubernetes controllers: the E2E Intent Orchestrator, the OCUDU Operator, and Nephio's free5GC Operator. The orchestrator translates \texttt{E2EQoSIntent} resources and coordinates the southbound paths. The two domain operators reconcile the RAN and core workloads as RAN NSSMF and Core NSSMF/VNFM components.

\subsection{Controller Lifecycle}

The controller follows the standard Kubernetes reconciliation pattern. When an \texttt{E2EQoSIntent} is created or updated, the reconciler fetches the current object, checks its generation, computes the desired translated state, and compares it with the last observed status. If the status already matches the current generation, the reconciliation ends without repeating southbound operations. Otherwise, the reconciler runs translation and domain realization and then writes a new status update.

This idempotent structure is important for GitOps integration. Porch and Config Sync may introduce asynchronous delay, and the Kubernetes API may deliver repeated events for the same resource version. The orchestrator therefore treats every reconciliation as a convergence attempt rather than a one-shot command. A failed stage records a reason and is retried according to the controller runtime backoff policy.

\begin{table}[t]
  \caption{Controller stages and observable artifacts}
  \label{tab:controller_stages}
  \begin{tabular}{lll}
    \toprule
    Stage & Main action & Observable artifact \\
    \midrule
    Admit & Check policy table & Status condition \\
    Translate & Derive core/RAN tuple & Intent status \\
    Mutate RAN & Patch KRM package & PackageRevision \\
    Configure core & Bind UE QoS data & Subscriber profile \\
    Observe & Read domain state & Fulfillment state \\
    \bottomrule
  \end{tabular}
\end{table}

Table~\ref{tab:controller_stages} summarizes the reconciliation pipeline. The implementation keeps these stages explicit in code rather than hiding them behind one monolithic handler. This helped testing because each stage can be exercised with a small fixture: an unsupported intent for admission, a known eMBB/URLLC input for translation, a package containing a cell configuration for RAN mutation, and a subscriber record for the core path. The same stage boundaries are also reflected in log messages and status reasons, which makes a failed run easier to inspect from \texttt{kubectl}.

\subsection{RAN Package Mutation}

For the RAN path, the orchestrator uses Nephio Porch as the package control point. The workflow has four steps. First, the orchestrator copies the blueprint package and pulls a mutable package revision. Second, it patches the structured KRM resource that represents OCUDU cell configuration. Third, it pushes and proposes the package revision. Fourth, it approves the package so that the GitOps pipeline can synchronize it to the workload cluster.

Listing~\ref{lst:ran_mutation} shows the resulting RAN slice configuration. The eMBB slice uses SD \texttt{0x010203}, a maximum PRB policy ratio of 50, and priority 10. The URLLC slice uses SD \texttt{0x112233}, a maximum PRB policy ratio of 100, and priority 200.

\begin{lstlisting}[caption={Mutated OCUDU slice configuration generated from intent translation}, label={lst:ran_mutation}]
spec:
  slicing:
    - sst: 1
      sd: 66051
      schedCfg:
        minPrbPolicyRatio: 0
        maxPrbPolicyRatio: 50
        priority: 10
    - sst: 1
      sd: 1122867
      schedCfg:
        minPrbPolicyRatio: 0
        maxPrbPolicyRatio: 100
        priority: 200
\end{lstlisting}

After synchronization, the OCUDU Operator renders the final \texttt{gnb-config.yml} and updates the gNB-related Kubernetes objects. Nephio and Config Sync deliver the KRM objects, and the operator performs the domain-specific realization needed to make that desired state executable by OCUDU. This preserves Git as the source of truth and allows package-level review or rollback.

\begin{figure}[t]
  \centering
  \includegraphics[width=\linewidth]{2.2.3-nephio-workflow-new.png}
  \caption{Nephio-style package workflow used by the RAN orchestration path}
  \Description{A package workflow from management cluster package mutation through Git synchronization to workload cluster application.}
  \label{fig:nephio_workflow}
\end{figure}

Figure~\ref{fig:nephio_workflow} illustrates why the RAN path is modeled as package mutation rather than direct pod modification. The package revision becomes the reviewable artifact, while the workload cluster only consumes approved state. This provides a clean operational boundary: the intent controller decides what RAN policy should change, and the GitOps pipeline decides when the approved configuration reaches the target cluster.

The mutation logic operates on structured YAML objects instead of template strings. The orchestrator locates the cell configuration resource, replaces only the \texttt{spec.slicing} entries, and preserves unrelated physical-layer fields such as band, ARFCN, bandwidth, and subcarrier spacing. It also updates the PLMN slice list so that the broadcast network identity remains consistent with the scheduler configuration. This is a small but important engineering detail: a scheduler entry alone is insufficient if UEs cannot discover or request the same S-NSSAI during registration and session establishment.

The resulting package is intentionally self-contained. A reviewer can inspect the package revision and see the old and new slice entries before approval, while the workload cluster receives the same artifact through Config Sync. This differs from a direct imperative update to a running ConfigMap, where the final state may be correct but the path that produced it is harder to audit or roll back.

\subsection{free5GC Operator and Core-Domain Configuration}

The Core NSSMF path uses the official free5GC Operator provided by Nephio. It watches \texttt{NFDeployment} resources, runs a reconciliation loop, and maintains the desired deployment state of free5GC network functions. This gives the core domain the same Kubernetes-native lifecycle behavior used by other Nephio-managed workloads.

Dynamic subscriber and QoS state follows a complementary path because it is not fully represented by the operator's deployment resources. The orchestrator uses the free5GC WebConsole API to configure S-NSSAI, subscriber binding, PDU session QoS, QFI, and 5QI for the target UEs. Separating NF deployment from subscriber configuration keeps the Core NSSMF role explicit without attributing unsupported runtime functions to the official operator.

The core implementation performs consistency checks before reporting success. For each target UE, it verifies that the intended S-NSSAI and QoS profile are present in the core-domain state. This does not prove that every packet will receive the desired treatment, because runtime QoS still depends on UPF and kernel behavior, but it does prove that the control-plane configuration requested by the intent is present.

The core path is organized around UE binding. For every \texttt{targetUEs} entry in an intent group, the orchestrator ensures that the subscriber profile contains the intended S-NSSAI and that the session management data contains the selected 5QI/QFI values. If the UE profile already exists, the orchestrator updates only the slice-related fields required by the active intent. If the profile is absent in the testbed, it can create the minimal subscriber state needed for the experiment. This behavior makes repeated experiments reproducible because the same intent converges to the same core-domain state.

The implementation treats this path as a compatibility layer rather than a permanent architectural preference. A fully declarative core package would be cleaner because RAN and core changes could share the same package lifecycle. However, forcing all free5GC subscriber data into an incomplete package abstraction would make the prototype less reliable. The hybrid path therefore preserves the most important property for this paper: a single high-level intent produces aligned core and RAN configuration, even though the two domains currently require different delivery mechanisms.

\subsection{OCUDU Operator and Hierarchical RAN Resources}

\texttt{OCUDUConfig} captures application-level settings such as CU/DU deployment and radio-emulation parameters. \texttt{OCUDUCellConfig} captures cell-level settings, including slice-aware scheduling. \texttt{PLMNConfig} captures network identity and S-NSSAI information. Splitting the configuration this way avoids large monolithic resources and lets the operator validate each layer independently before producing the final ConfigMap.

The OCUDU Operator performs reference walking across these resources. Starting from the deployment resource, it resolves the referenced application, cell, and PLMN objects, validates that the referenced slices are mutually consistent, and then renders the final gNB configuration. When rendering succeeds, it updates Kubernetes resources and marks the deployment as ready. When a referenced object is missing or inconsistent, it records a condition instead of producing a partial configuration. This local check is separate from the E2E orchestrator's admission policy and protects the workload cluster from broken package content.

\begin{lstlisting}[caption={Representative rendered scheduler fragment}, label={lst:rendered_scheduler}]
cell_cfg:
  slicing:
    - sst: 1
      sd: 0x010203
      sched_cfg:
        min_prb_policy_ratio: 0
        max_prb_policy_ratio: 50
        priority: 10
    - sst: 1
      sd: 0x112233
      sched_cfg:
        min_prb_policy_ratio: 0
        max_prb_policy_ratio: 100
        priority: 200
\end{lstlisting}

Listing~\ref{lst:rendered_scheduler} shows the final form consumed by the RAN software. Comparing this artifact with the original intent and the intermediate package gives a concrete end-to-end trace: service expectation, translated profile, KRM package, and rendered runtime configuration.

\subsection{Failure Handling}

The prototype distinguishes three failure classes. Translation failures occur before any southbound operation, such as an unsupported slice type or an impossible combination of expectations. Delivery failures occur when a package cannot be mutated, approved, or synchronized. Domain failures occur when either domain operator cannot reconcile its delivered resources. These classes are written into intent status so that the consumer can tell whether the problem is semantic, pipeline-related, or domain-local.

This failure taxonomy is deliberately aligned with future 3GPP intent features. Feasibility checks and conflict reports can be added before translation, while negotiation and utility functions can be added as policy modules that select among alternative translations. The current implementation is therefore a first step toward richer intent handling rather than a closed one-off script.
